\scanpage{248}
En termes de $\sigma\in H$, \unsure{cel|cela} s'exprime par

\[
\left\{\begin{aligned}
&\sigma\text{ normalise }T_H,\text{ et }\deleted{\mathrm{int}(\sigma)\mid T_H/\unsure{\mathfrak z|\zeta}}\ \forall t\in T_H,\text{ on a}\\
&\qquad\sigma(t)\equiv t^{-1}\quad(\mathfrak z_H)\\
&\text{donc }\sigma(t)=t^{-1}\quad\text{si }\mathfrak z_H=\{1\}\quad.
\end{aligned}\right.
\tag*{(86)}
\]

\[
B_{1H}=\sigma(B_{0H}),\quad B_{0H}=\sigma(B_{1H})\quad.
\tag*{(87)}
\]

Notons maintenant

(88) La donnée, \deleted{de} pour un graphe homogène défini par $(H,U_0\subset H,\sigma\in H)$ (satisfaisant $\sigma^2\in U_0\cap\sigma(U_0)=\mathfrak z$) d'un groupe d'homothéties $Z$, \inserted{satisfaisant à (72)\textsuperscript{*}} \uline{équivaut} à celle d'un sous-groupe $T_H\subset H$, normalisant $U_0$, tel qu'on ait $T_H\cap U_0=\mathfrak z_H$, $[$que $T_H/\mathfrak z_H$ soit commutatif$]$, \deleted{et} que $\sigma$ normalise $T_H$ et \deleted{\unsure{ind|tant}} satisfasse $\sigma(t)\equiv t^{-1}\ \deleted{\unsure{\mathrm{mod}\,\mathfrak z}}\ \inserted{\mathrm{mod}\,\mathfrak z_H}$, pour $t\in T_H$, et tel enfin qu'on ait $\underbrace{TU_0}_{B_{0H}}\cap U_1=\mathfrak z_H$ (d'où \deleted{\unsure{par}} \inserted{en transformant par $\sigma$} $TU_{\unsure{1}}\cap U_0=\mathfrak z_H$)

% slanted
\margin{* i.e. avec opération \uline{fidèle} de $Z$, sur l'ens $\pi$ des structures de graphe engendré par $(S,\overrightarrow{A_0})$ sous $Z$}

% slanted
\margin{$[\ ]$ résulte du reste}

% slanted
\margin{\uline{NB} On voit alors que \deleted{\unsure{(89)}}}

En effet, si $T_H$ est \deleted{un} sous groupe de $H$ normalisant $U_0$, alors $\inserted{Z\overset{\mathrm{def}}{=}}T_H/U_0\cap T_H\simeq\underbrace{T_H\mathbin{.}U_0}_{B_{0H}}/U_0$, opère \inserted{fidèlement} sur l'espace homogène $S\simeq\deleted{\unsure{G|O}}H/U_0$, il faut vérifier que c'est un groupe d'homothéties, on a \inserted{supposé} la commutativité, d'autre part pour $t\in T_{N}$ on a $\sigma(t)\mathbin{.}t\in\mathfrak z_H\subset U_{1}\cap U_0\inserted{\subset U_1\mathbin{.}U_0}$, un calcul déjà fait, et immédiat nous montre que \unsure{elle|celle|cela} implique que $t$ est une homothétie. \deleted{A}

\deleted{\unsure{Le cas vu|le cas vu} \unsure{alors}, on a utilisé seulement l'inclusion $T_H\cap U_0\subset\mathfrak z$,}

% archive sheet number: 247
